% Generated draft: requires mathematical and source review.
\documentclass[11pt]{article}
\usepackage[margin=25mm]{geometry}
\usepackage{fontspec}
\setmainfont{Latin Modern Roman}
\usepackage{amsmath,amssymb,mathtools,unicode-math}
\setmathfont{Latin Modern Math}
\usepackage{enumitem,needspace,xurl,hyperref}
\setlength{\parindent}{0pt}
\setlength{\parskip}{5pt}
\setlength{\emergencystretch}{3em}
\newcommand{\N}{\mathbb{N}}
\newcommand{\Z}{\mathbb{Z}}
\newcommand{\Q}{\mathbb{Q}}
\newcommand{\R}{\mathbb{R}}
\newcommand{\C}{\mathbb{C}}
\newcommand{\F}{\mathbb{F}}
\newcommand{\A}{\mathbb{A}}
\newcommand{\PP}{\mathbb{P}}
\newcommand{\E}{\mathbb{E}}
\newcommand{\eps}{\varepsilon}
\newcommand{\dd}{\,\mathrm{d}}
\begin{document}
\section*{UnsolvedMath ID: 30002200}
\textbf{Source problem code:} catalogue entry
\subsection{Definitions and mathematical statement}
Interpret the torus rank as the dimension of an acting torus \(T\cong(S^1)^r\). Its rational equivariant coefficient ring is
\[
R_r=H^*(BT;\Q)\cong\Q[t_1,\ldots,t_r],\qquad |t_i|=2.
\]
For a \(T\)-space \(X\), its rational Borel equivariant cohomology is the graded \(R_r\)-module
\[
M_X=H_T^*(X;\Q):=H^*(ET\times_T X;\Q),
\]
where \(ET\) is a contractible space with a free \(T\)-action and \(BT=ET/T\).

Under the usual connected convention, a rational Poincar\'e duality space of formal dimension \(n\geq0\) has finite-dimensional total rational cohomology, \(H^0(X;\Q)\cong H^n(X;\Q)\cong\Q\), no cohomology above degree \(n\), and perfect cup-product pairings
\[
H^i(X;\Q)\otimes H^{n-i}(X;\Q)\longrightarrow H^n(X;\Q)\cong\Q
\qquad(0\leq i\leq n).
\]
The supplied record does not specify whether this connected convention or a broader convention is intended.

For a finitely generated graded \(R_r\)-module \(M\), call \(M\) a \(0\)-th syzygy automatically. For \(k\geq1\), call it a \(k\)-th syzygy if there is an exact sequence of graded modules
\[
0\longrightarrow M\longrightarrow F_{k-1}\longrightarrow\cdots\longrightarrow F_0\longrightarrow N\longrightarrow0,
\]
where the \(F_i\) are finite free graded \(R_r\)-modules, allowing grading shifts, and \(N\) is finitely generated. Set
\[
\operatorname{syzord}_{R_r}(M)=\sup\{k\in\Z_{\geq0}:M\text{ is a }k\text{-th syzygy}\}.
\]

The supplied record asserts that, for the rational Poincar\'e duality spaces and torus actions intended by its source,
\[
M_X\text{ not free over }R_r
\quad\Longrightarrow\quad
\operatorname{syzord}_{R_r}(M_X)\leq
\left\lfloor\frac{r-1}{2}\right\rfloor.
\]
The natural rank-by-rank interpretation of its sharpness question is: for every integer \(r\geq5\), does there exist an admissible rational Poincar\'e duality space \(X\) with an action of \(T\cong(S^1)^r\) such that \(M_X\) is finitely generated, is not free over \(R_r\), and satisfies
\[
\operatorname{syzord}_{R_r}(M_X)=\left\lfloor\frac{r-1}{2}\right\rfloor?
\]
This is a conditional formulation: the record does not specify the admissible category of spaces and actions, the relevant finiteness hypotheses, or whether effectiveness is required. It also does not disambiguate acting-torus dimension from other meanings of torus rank. These source hypotheses must be checked before the displayed question is treated as an exact transcription.
\subsection{Short English statement}
The supplied question asks whether the stated upper bound on the syzygy order of nonfree rational equivariant cohomology can be attained for every torus rank \(r\geq5\). The formulation uses the acting-torus and finitely generated module conventions described above; the precise admissible spaces and actions remain unspecified.
\subsection{Sources}
\begin{itemize}
\item Supplied local UnsolvedMath registry record, permanent ID \(30002200\), titled ``Sharp Syzygy Bounds for Torus Actions,'' in the Oberwolfach Reports Open Problems collection; exact source problem code: \texttt{catalogue entry}. No individual proposer is named.
\item Supplied source citation: \emph{Mini-Workshop: Cohomology Rings and Fundamental Groups of Hyperplane Arrangements, Wonderful Compactifications, and Real Toric Varieties} (2013), \emph{Oberwolfach Reports}, pp.\ 2939--2983, DOI \texttt{10.4171/owr/2012/49}. Supplied source URL and literature-source link, labeled ``original OWR report'': \url{https://doi.org/10.4171/owr/2012/49}. The record gives proposed year 2012; its citation labels the report 2013.
\item Literature assessment dated exactly \texttt{2026-08-21}: the supplied curation labels the problem open, reports sharpness for \(r\leq4\) and an unresolved question for \(r\geq5\), and states that no later complete construction or proof was located. This is supplied dated curation, not an independently verified literature assessment. The background describes it as literature triage requiring independent verification of time-sensitive status and cited claims.
\item The registry assigns the statement status \texttt{corrected\_verified} and states that its mathematical content was checked against the official report, with notation, context, and question form normalized without changing the source-level claim. These are supplied verification assertions; the report has not been independently consulted for this draft.
\end{itemize}
\end{document}
